\subsection{I-32. Exponential coordinate-cone compression is necessary even at the common center}
\label{I-32}

\paragraph{Statement.}
For every $w\geq3$ and every coefficient field, put
\[
 h=3\cdot2^{w-2}-2,\qquad n=2h.
\]
There is an explicit width-$w$ word of $n$ univariate linear
matrices in distinct variables $x_1,y_1,\ldots,x_h,y_h$ with
the following properties. Image cones use the algebraic-closure
convention of I-31.

Every nonempty internal interval has generic rank in
$\{2,\ldots,w\}$, and that rank is preserved at every common
center $x_i=y_i=\tau$. In particular it is preserved at the
transcendental center $t\mathbf1$ used in I-31. Nevertheless,
after translation to residual variables at that center, a
restriction obtained by setting unselected coordinates to zero
preserves the full matrix image cone if and only if it retains
at least one coordinate from each pair $(x_i,y_i)$. The minimum
number of retained positions is therefore exactly $h$.

Thus a uniform polynomial-in-$w$ coordinate-cone compression
bound is false, even at a common transcendental center where
all internal interval ranks are preserved. In terms of the
number $L=w-1$ of allowed rank levels, the example needs
$3\cdot2^{L-1}-2$ positions. It precludes replacing the rank-level
recursion by a polynomial bound for this exact compression
criterion. It does not preclude changing the center or using
other dense maps: the distinguished nonzero scalar below is
detected by one centered SV pair.

\paragraph{1. The seed word.}
We recall the explicit matrices of I-24, including the property
needed here. Set $P_r=\operatorname{diag}(I_r,0_{w-r})$ and
\[
 Q_{r,a}(s,u)=(P_r+sE_{w,a})(P_r+uE_{a,w})
       \quad(2\leq r<w,\ a\leq r).
\]
With fresh variables at every occurrence, define
\[
 H_w=\text{empty word},\qquad
 H_k=H_{k+1}Q_{k,k}H_{k+1}\quad(3\leq k<w),
\]
and
\[
 W=H_3Q_{2,1}H_3Q_{2,2}H_3.
\]
If $L_k$ is the length of $H_k$, then
$L_w=0$, $L_k=2L_{k+1}+2$, and
$3L_3+4=h$. Label the successive variables of this word by
$Y_1,\ldots,Y_h$. Each factor is of one of the two forms
\[
 P_r+Y_iE_{w,a},\qquad P_r+Y_iE_{a,w}.
 \tag{I.54}
\]
They have generic rank $r$ and constant matrix $P_r$.

For completeness, if a row $q$ has $q_a=0$, then
\[
 qQ_{r,a}(s,u)
   =\sum_{j\leq r}q_j e_j^T
       +q_ws e_a^T+q_wsu e_w^T.
\]
Starting with coordinates $k,\ldots,w-1$ zero, $H_k$ leaves
the earlier coordinates unchanged and multiplies the last
coordinate by all its variables. This follows by induction:
the first $H_{k+1}$ leaves coordinate $k$ zero, $Q_{k,k}$
performs the displayed transfer and resets all higher active
coordinates, and the second $H_{k+1}$ applies again. In the
full word starting at $e_w^T$, the two rank-two gadgets use
successively the still-zero coordinates one and two and reset
the coordinates required by the intervening $H_3$ words.
Consequently
\[
 e_w^T W(Y_1,\ldots,Y_h)e_w=\prod_{i=1}^hY_i.
 \tag{I.55}
\]
This calculation is valid in every characteristic.

\paragraph{2. A paired lift with preserved common-center ranks.}
Replace every factor in (I.54) by two factors on fresh variables.
For $E=E_{w,a}$, use
\[
 (I_w+x_iE)\bigl((I_w-y_iE)P_r\bigr),
 \tag{I.56}
\]
and for $E=E_{a,w}$, use
\[
 \bigl(P_r(I_w+x_iE)\bigr)(I_w-y_iE).
 \tag{I.57}
\]
Since $E^2=0$, all matrices $I_w+zE$ are invertible with
inverse $I_w-zE$. In (I.56), $EP_r=E$; in (I.57),
$P_rE=E$. Each pair therefore equals
$P_r+(x_i-y_i)E$. The full lifted word $\widetilde W$ obeys
the matrix identity
\[
 \widetilde W(x_1,y_1,\ldots,x_h,y_h)
       =W(x_1-y_1,\ldots,x_h-y_h).
 \tag{I.58}
\]
Every factor is still univariate and linear. Exactly one half
of each pair is invertible; the other is an invertible matrix
times $P_r$, on the appropriate side, and has rank exactly $r$
for every parameter value.

We verify every partial interval, not only complete pairs.
At $x_i=y_i=\tau$, each complete pair is $P_r$.
A partial initial pair can only contribute its second half.
For (I.56) that half is $(I_w-\tau E)P_r$; for (I.57) it is
$I_w-\tau E$. A partial final pair can only contribute its
first half. For (I.56) it is $I_w+\tau E$; for (I.57) it is
$P_r(I_w+\tau E)$. All intervening pairs are complete.
Thus the constant interval product has the form
\[
 U_{\rm left}\left(\prod P_{r_i}\right)U_{\rm right},
 \tag{I.59}
\]
where the two boundary matrices are invertible, and the
projections are exactly those attached to singular halves
lying in the interval. The same description handles an
interval contained in one pair: a single half has its
displayed rank, and the complete pair is $P_r$. If no singular
half occurs, the interval is a product of invertible matrices.

The nested diagonal projections in (I.59) have product
$P_q$, where $q$ is the minimum rank of the singular halves
present; take $q=w$ if there are none. The common-center rank
is therefore exactly $q$. The generic interval rank is at
most $q$, because every singular half has that fixed rank
and products cannot increase it. It is at least $q$ because
the specialization (I.59) has rank $q$, so some $q$-minor is
a nonzero polynomial before specialization. These two bounds
prove equality. This argument applies to every $\tau$, in
particular to $\tau=t$ over the coefficient field extended
by a transcendental element.

\paragraph{3. The exact coordinate-compression threshold.}
Write $x_i=t+X_i$, $y_i=t+Y_i$. Formula (I.58) gives
\[
 \widetilde W(t+X_1,t+Y_1,\ldots,t+X_h,t+Y_h)
     =W(X_1-Y_1,\ldots,X_h-Y_h),
\]
and (I.55) becomes
\[
 e_w^T\widetilde W(t+\mathbf X,t+\mathbf Y)e_w
       =\prod_{i=1}^h(X_i-Y_i)\ne0.
 \tag{I.60}
\]
Let a coordinate restriction retain a set of residual positions
and set all other residual coordinates to zero. If it omits
both positions of any pair, the displayed scalar is identically
zero after restriction. The restricted matrix image cone is
then contained in the homogeneous hyperplane
$\{M:e_w^TMe_w=0\}$. The original cone is not contained in
that hyperplane by (I.60). Hence the two cones differ.

Conversely, if at least one position in every pair is retained,
the remaining differences $X_i-Y_i$ can still be assigned
independently and arbitrarily. Formula (I.58) shows that the
restricted full matrix image equals the original full matrix
image, not merely its cone closure. This proves the exact
criterion and the minimum $h$: retaining $X_1,\ldots,X_h$
alone attains it.

The counterexample lies in the actual final rank-one test
class. Append $L=e_1e_w^T$ and $R=e_we_1^T$ to the word.
Its full product is (I.60) times $e_1e_1^T$ and has generic
rank one. Every prefix beginning at $L$ is nonzero, since
otherwise the full product would vanish, and therefore has
rank one. Every interval wholly inside $\widetilde W$ has
rank at least two. Thus the shortest rank-one suffix of the
prefix just before $R$ begins at $L$, and adjoining $R$
gives the full word as a local final test. Its entire internal
higher-rank profile is already preserved at the common center.

\paragraph{4. One dense pair still detects the scalar.}
Choose distinct interpolation nodes for the $2h$ original
variables, and denote their Lagrange polynomials by
$\lambda_{x_i}$ and $\lambda_{y_i}$. The common-center
one-pair substitution gives
\[
 \prod_{i=1}^h(x_i-y_i)
   \longmapsto
 z^h\prod_{i=1}^h
       \bigl(\lambda_{x_i}(y)-\lambda_{y_i}(y)\bigr)\ne0.
 \tag{I.61}
\]
Each difference is nonzero because its value at the node
for $x_i$ is one. Their product is nonzero in the parameter
polynomial ring over every working field with these distinct
nodes. Numerical evaluation uses an adequate grid or extension,
as in I-21; no assertion is made that every node value detects
this product.

This separates the refuted exact coordinate-deletion criterion
from dense nonvanishing. In particular it rules out a polynomial
dimension charge for uniformly preserving the full matrix cone
by retaining original variable positions at every preserved
common center. It does not rule out a smaller bound after a
different dense shift, a weaker sufficient invariant than full
cone equality, or a more efficient sampling argument. The
constant-center difference cancellation is explicit, not an
assertion about every possible preserved center. $\square$

\paragraph{Derivation record.}
The paired lift and the common-center cone obstruction were
derived by the purity track from the independently proved
I-24 matrices. The result is a deduction within this report;
no literature-priority claim is made.
